\documentclass{exam}

\newif\ifA
\newif\ifkey
\newif\ifprintselected   % required by exam_config's \ansbox

%%%%%%%%%%%%%%%%%%%%%%%%%%%%%%%%%%%%%%%%%%%%%%%%%%%%%%
%%% SET THESE IF VARIABLES
%%%%%%%%%%%%%%%%%%%%%%%%%%%%%%%%%%%%%%%%%%%%%%%%%%%%%%
% Lab Num -- Module 3's six lab periods follow L04 (Lab 1) and the two
% Module 2 measurement labs, so L17 is Lab 4. Override with \def\labNumIn{}
% if the course sequence renumbers.
\ifdefined\labNumIn
	\def\labNum{\labNumIn}
\else
	\def\labNum{4}
\fi

% is this the key or not?
\ifdefined\iskey
	\ifnum\iskey=1
		\keytrue
	\else
		\keyfalse
	\fi
\else
	%% Pick one manually
	\keyfalse % if this is NOT the key
%	\keytrue     % if this IS the key
\fi

% set up the header and footer - change for every term
\ifdefined\thisTermIn
	\def\thisterm{\thisTermIn}
\else
	\def\thisterm{Fall 2026}
\fi
%
\def\thisclass{ECE 444}

%%%%%%%%%%%%%%%%%%%%%%%%%%%%%%%%%%%%%%%%%%%%%%%%%%%%%%
%%% END SET THESE IF VARIABLES
%%%%%%%%%%%%%%%%%%%%%%%%%%%%%%%%%%%%%%%%%%%%%%%%%%%%%%
% packages and shortcuts
\input{myPackages_gr}
\usepackage[hidelinks, linkcolor=red]{hyperref}
\input{myShortcuts}
\setlength{\parindent}{0pt}
\setlength{\parskip}{6pt}

%% exam class customization
\input{exam_config}
\printselectedfalse      % full worked solutions in the key, not answers-only
\CorrectChoiceEmphasis{\color{red}}
\SolutionEmphasis{\color{red}}
\renewcommand\questionlabel{\textbf{\thequestion.}}
\renewcommand{\thepartno}{\alph{partno}}
\renewcommand{\partlabel}{(\thepartno)}

% Fillable table cell: blank in the student copy, the expected value in red on
% the KEY. Fixed-width centered columns keep the blank cells writable.
\newcolumntype{K}[1]{>{\centering\arraybackslash}p{#1}}
\newcommand{\kv}[1]{\ifkey{\color{red}#1}\fi\rule[-0.10in]{0pt}{0.34in}}

% print the answers if this is the key
\ifkey
	\printanswers % if this IS the key
\else
	\noprintanswers  %if this is NOT the key
\fi

% print the header and footer
% need to print key on top if this is the key
\ifkey
	\firstpageheader{}{\color{red}{ KEY}}{}
	\runningheader{\thisclass\hspace{1pt} -- \thisterm }{Lab \#\labNum\hspace{1pt}\color{red}{ KEY}}{\thepage\hspace{1pt} of\hspace{1pt} \numpages}
	\firstpagefooter{}{}{}
	\runningfooter{}{}{}
\else
	\firstpageheader{}{}{}
	\runningheader{\thisclass \hspace{1pt} -- \thisterm }{Lab \#\labNum\hspace{1pt}}{\thepage\hspace{1pt} of\hspace{1pt} \numpages}
	\firstpagefooter{}{}{}
	\runningfooter{}{}{}
\fi

\runningheadrule

\begin{document}
\pagestyle{headandfoot}   % exam class

%%%%%%%%%%%%%%%%%%%%%%%%%%%%%%%%%%%%%%%%%%%%%%%%%%%%%%
% title block
%%%%%%%%%%%%%%%%%%%%%%%%%%%%%%%%%%%%%%%%%%%%%%%%%%%%%%

\begin{center}
USAF ACADEMY \\
DEPARTMENT OF ELECTRICAL AND COMPUTER ENGINEERING \\
LAB \#\labNum\ (Lesson 17) - Introduction to Phased Array Hardware \\
\textbf{Lab Sheet} \\
\thisterm \\[4mm]
\end{center}

\vspace{-4pt}
\textbf{Name:} \rule[-2pt]{2.6in}{0.4pt} \hfill
\textbf{Section:} \rule[-2pt]{0.9in}{0.4pt} \hfill
\textbf{Date:} \rule[-2pt]{1.1in}{0.4pt}

\vspace{4pt}

\fullwidth{\textbf{Learning Objective 3.3: } I can identify the hardware architecture
of the ADALM-PHASER and control it via SDR software.}

\vspace{2pt}

This sheet is the turn-in document for the Lesson 17 lab. Fill in every box, and
show the arithmetic where a box asks for a calculated value. The lab is a team
assignment, so one sheet per team.

\begin{questions}

%%%%%%%%%%%%%%%%%%%%%%%%%%%%%%%%%%%%%%%%%%%%%%%%%%%
\question \textbf{The frequency plan.} Record the frequency the GUI reports for
your HB100 after pressing \textbf{Find HB100}, then work the high-side mixing
arithmetic that follows from it. The IF is fixed by the filtering after the
mixer, so the LO is the only frequency that moves when the source moves.

\begin{center}
\renewcommand{\arraystretch}{1.25}
\begin{tabular}{| l | K{2.4in} |}
\hline
\textbf{Quantity} & \textbf{Value} \\\hline\hline
Measured HB100 frequency $f_{\text{RF}}$ & \kv{your own, in $10.1$--$10.7\ \ghz$} \\\hline
Required LO frequency $f_{\text{LO}}$ & \kv{$f_{\text{RF}} + 2.2\ \ghz$} \\\hline
Resulting IF $f_{\text{IF}}$ & \kv{$2.200\ \ghz$} \\\hline
Is the LO inside $12.2$--$13.0\ \ghz$? & \kv{yes, for any HB100 in that band} \\\hline
\end{tabular}
\end{center}

Show the arithmetic for your measured frequency here.

\begin{solutionorlines}[1.1in]
For a nominal source at $10.525\ \ghz$, high-side injection puts the LO at
\eq{
	f_{\text{LO}} &= f_{\text{RF}} + f_{\text{IF}} = 10.525 + 2.200 = 12.725\ \ghz \\
	f_{\text{IF}} &= f_{\text{LO}} - f_{\text{RF}} = 12.725 - 10.525 = 2.200\ \ghz
}
That LO sits inside the $12.2$ to $13.0\ \ghz$ VCO range, so the setting is
reachable. Any measured frequency in $10.1$ to $10.7\ \ghz$ gives an LO inside
the range, which is why the hardware covers every HB100 in the drawer.
\end{solutionorlines}

%%%%%%%%%%%%%%%%%%%%%%%%%%%%%%%%%%%%%%%%%%%%%%%%%%%
\newpage
\question \textbf{The labeled block diagram.} Sketch the receive chain from the
patches to the Raspberry Pi. Label every block with its name and with the
frequency present at that point, mark where the analog summing happens, and mark
where the two digital channels begin.

\begin{solutionorbox}[3.4in]
Patches ($f_{\text{RF}}$, eight of them, $d = 14$ mm apart) $\rightarrow$ ADL8107
LNA behind each patch ($f_{\text{RF}}$) $\rightarrow$ two ADAR1000 beamformers,
four elements each, phase and gain per element, \textbf{analog summing happens
here} ($f_{\text{RF}}$, two outputs) $\rightarrow$ two LTC5548 mixers driven by
the ADF4159 PLL and HMC735 VCO at $f_{\text{LO}} = f_{\text{RF}} + 2.2\ \ghz$
($12.2$ to $13.0\ \ghz$) $\rightarrow$ two IF channels at $2.2\ \ghz$
$\rightarrow$ ADALM-Pluto (AD9361), tuned to $2.2\ \ghz$, sampling at 3 MSPS,
\textbf{the digital channels begin here} $\rightarrow$ Raspberry Pi over USB,
running the backend and serving the browser interface.
\end{solutionorbox}

%%%%%%%%%%%%%%%%%%%%%%%%%%%%%%%%%%%%%%%%%%%%%%%%%%%
\newpage
\question \textbf{FFT observations.} Work on the FFT tab with the source at
$1\ \m$ and record what the spectrum does as you change two controls.

\begin{parts}

	\part Where does the tone sit in the baseband window, and how far above the
	noise floor is it at the preset Rx Gain?

	\vspace{6pt}
	\begin{center}
	Peak offset = \ansbox[1.7in]{$\approx 1\ \mhz$} %
	Peak above floor = \ansbox[1.7in]{$\ge 20\db$}
	\end{center}
	\vspace{6pt}

	\part Step \textbf{Rx Gain} down 10 dB and up 10 dB from its preset value and
	record all three settings. Report the peak and floor levels in dBFS and the
	separation between them in dB.

	\vspace{4pt}
	\begin{center}
	\renewcommand{\arraystretch}{1.25}
	\begin{tabular}{| l | K{1.25in} | K{1.25in} | K{1.35in} |}
	\hline
	\textbf{Rx Gain} & \textbf{Peak (dBFS)} & \textbf{Floor (dBFS)} & \textbf{Separation (dB)} \\\hline\hline
	Preset $-10$ dB & \kv{$10$ dB below the reference row} & \kv{$10$ dB below the reference row} & \kv{1--2 dB less than the reference row} \\\hline
	Preset & \kv{reference} & \kv{reference} & \kv{$\ge 20$ dB} \\\hline
	Preset $+10$ dB & \kv{$10$ dB above the reference row} & \kv{$10$ dB above the reference row} & \kv{$\ge 20$ dB} \\\hline
	\end{tabular}
	\end{center}
	\vspace{4pt}

	\begin{solution}
	Both levels move together with the slider, because this gain is applied
	inside the SDR, long after the LNAs have set how much noise rides on the
	signal. The separation therefore holds within a decibel or two across the
	three settings. It shrinks at the lowest setting, where the SDR's own noise
	starts to contribute. Any set of readings whose separation changes by more
	than a few dB across the three rows points to a saturating or misaimed setup.
	\end{solution}

	\part Change \textbf{Signal Freq (GHz)} by $-0.0005\ \ghz$ and record how far
	the peak moves and in which direction.

	\vspace{6pt}
	\begin{center}
	Peak shift = \ansbox[2.0in]{$-500\ \khz$}
	\end{center}
	\vspace{6pt}

	\begin{solutionorlines}[0.9in]
	The GUI holds the IF fixed and moves the LO, so lowering the assumed source
	frequency by 500 kHz lowers the LO by the same amount and the tone appears
	500 kHz lower in the baseband window. The window is only 3 MHz wide at a
	3 MSPS sample rate, so an error larger than about 1.5 MHz walks the tone off
	the display, which is what a bad \textbf{Find HB100} result looks like.
	\end{solutionorlines}

\end{parts}

%%%%%%%%%%%%%%%%%%%%%%%%%%%%%%%%%%%%%%%%%%%%%%%%%%%
\newpage
\question \textbf{Written answers.} Three or four sentences each.

\begin{parts}

	\part Why does the control software have to hunt for the HB100's frequency
	instead of assuming $10.525\ \ghz$?

	\begin{solutionorlines}[1.7in]
	The HB100's oscillator is a free-running dielectric resonator, not a locked
	synthesizer, so its output frequency is set by the mechanical dimensions of
	that resonator and drifts with temperature. Units land anywhere from about
	$10.1$ to $10.7\ \ghz$. Every LO setting and every steering phase the backend
	computes depends on the actual frequency, so the software sweeps the LO,
	watches where the IF tone appears, and records the answer instead of assuming
	one.
	\end{solutionorlines}

	\part The array has eight elements, but software sees only two digital
	channels. Explain where the other six went, and name one measurement this
	rules out.

	\begin{solutionorlines}[1.9in]
	Each ADAR1000 applies a phase and a gain to its four elements at RF and then
	sums those four into a single output, so the individual element signals stop
	existing before anything is digitized. Two subarray sums reach the Pluto, and
	the eight per-element signals cannot be recovered from them. That rules out
	any measurement that needs per-element data after the fact, such as forming
	several independent beams from one set of samples, or running an adaptive
	beamformer with more than one null. The MVDR run in Lesson 28 has exactly two
	digital degrees of freedom for this reason.
	\end{solutionorlines}

\end{parts}

\end{questions}

\vspace{1.0em}
\noindent\textbf{Documentation:} \rule[-2pt]{0.6\linewidth}{0.4pt}

\end{document}
